\documentclass[12pt,a4paper]{article}
\usepackage{amsmath,amssymb,amsthm}
\usepackage{bm}               % optional, for bold symbols
\usepackage{hyperref}
\usepackage{geometry}
\geometry{margin=1in}
\usepackage{array}
\usepackage{booktabs}
\usepackage{longtable}
\usepackage{enumitem}

%-------------------------------------------------
%  Macros
%-------------------------------------------------
\newcommand{\dfix}{\mathfrak{0}}          % 0d – the unique fixed point
\newcommand{\dist}{\|\cdot\|}            % distance symbol
\newcommand{\Rep}{\mathcal{R}}           % Representable objects
\newcommand{\Prf}{\mathcal{P}}           % Proof objects
\newcommand{\Tset}{\mathcal{T}}          % set of time tokens

%-------------------------------------------------
%  Theorem environments
%-------------------------------------------------
\newtheorem{definition}{Definition}[section]
\newtheorem{theorem}{Theorem}[section]

%-------------------------------------------------
\begin{document}

\title{Formalisation of ``Proof’’ via Monomorphisms, Epimorphisms and
Contractions}
\author{Masaomi Chiba}
\date{2026-02-02}
\maketitle

\begin{abstract}
We show that the act of proving can be described rigorously by
\begin{itemize}[nosep]
  \item a monomorphism (injective embedding) $\iota\colon\Rep\hookrightarrow\Prf$,
  \item an epimorphism (surjective projection) $\phi\colon\Prf\twoheadrightarrow\Rep$,
  \item and a contraction property of $\phi$ with respect to the distance
        $\|X\|=K(X)+\ell(X)$.
\end{itemize}
Together these give a categorical picture of proof, guarantee loss‑free
information transfer, and ensure monotone decrease of the distance until the
unique $0$‑dimensional fixed point $\dfix$ is reached.
We use this categorical core as the convergence mechanism behind the
\emph{Catuskoti} (tetralemma) interpretation: distinct truth statuses are
understood as collapsing, via evaluation, to the same $0$‑dimensional fixed
point (``0d'') $\dfix$.
In particular, we do not treat ``emptiness'' as a separate object: it is
simply this $0$‑dimensional fixed point.
\end{abstract}

%%%%%%%%%%%%%%%%%%%%%%%%%%%%%%%%%%%%%%%%%%%%%%%%%%%%%%%%%%%%
\section{What becomes more rigorous}
%%%%%%%%%%%%%%%%%%%%%%%%%%%%%%%%%%%%%%%%%%%%%%%%%%%%%%%%%%%%

\begin{definition}[Objects]
$$
\begin{aligned}
\Rep &:= \{\, (s_x,s_y)\mid s_x,s_y\in\Sigma \,\}
      &&\text{(static, time‑free information)}\\[2mm]
\Prf &:= \{\, ((s_x,t_x),(s_y,t_y))
          \mid (s_x,s_y)\in\Rep,\;t_x,t_y\in\Tset \,\}
      &&\text{(proof objects, time‑axis of length $1$)} .
\end{aligned}
$$
Here $\Sigma$ is a set of bit‑strings (natural numbers, symbols, …) and
$\Tset$ is a set of \emph{time tokens} (any objects with $\ell(t)=1$).
\end{definition}

\begin{definition}[Embedding and projection]
$$
\begin{aligned}
\iota &: \Rep\hookrightarrow\Prf,
        &\iota(s_x,s_y)&:=\bigl((s_x,t_x),(s_y,t_y)\bigr),
        &&t_x,t_y\in\Tset\text{ arbitrary};\\[2mm]
\phi &: \Prf\twoheadrightarrow\Rep,
        &\phi\bigl((s_x,t_x),(s_y,t_y)\bigr)&:=(s_x,s_y).
\end{aligned}
$$
\end{definition}

\begin{definition}[Distance (axiomatised)]
We fix two maps
\[
  K:\Rep\cup\Prf\to\mathbb{N},
  \qquad
  \ell:\Rep\cup\Prf\to\{0,1\}
\]
and define the distance
\[
  \|X\|\;:=\;K(X)+\ell(X).
\]
We require the following conventions.
\begin{enumerate}[label=(D\arabic*)]
  \item $\ell(x)=0$ for all $x\in\Rep$ and $\ell(p)=1$ for all $p\in\Prf$.
  \item (Time erasure does not change $K$.) For all $p\in\Prf$,
        \[
          K\bigl(\phi(p)\bigr)=K(p).
        \]
\end{enumerate}
In this note, $K$ is treated as an abstract \emph{description length} rather
than the full Kolmogorov complexity (so the above identity is admissible as a
modelling assumption).
\end{definition}

\begin{definition}[Monoid‑type properties]
\begin{itemize}[nosep]
  \item $\iota$ is a \textbf{monomorphism} (left‑cancellable injection);
        consequently no information is lost:
        $\iota(x)=\iota(y)\Rightarrow x=y$.
  \item $\phi$ is an \textbf{epimorphism} (right‑cancellable surjection);
        every static object is the image of some proof.
  \item $\iota$ is a \textbf{split monomorphism}:
        there exists $\sigma:\Prf\to\Rep$ with $\sigma\circ\iota=\operatorname{id}_{\Rep}$.
        One can take $\sigma((s_x,t_x),(s_y,t_y))=(s_x,s_y)$.
  \item $\phi$ is a \textbf{split epimorphism}:
        there exists $\rho:\Rep\to\Prf$ with $\phi\circ\rho=\operatorname{id}_{\Rep}$.
        Take $\rho=\iota$.
  \item $\iota\dashv\phi$ (they form an adjoint pair); $\iota$ adds a unit of time,
        $\phi$ removes it.
\end{itemize}
\end{definition}

\begin{definition}[Contraction (time erasure)]
For every $p\in\Prf$,
\[
  \|\phi(p)\| = \|p\|-1.
\]
This is immediate from (D1)--(D2): $\phi$ removes exactly one unit of
``time length'' while preserving $K$.
\end{definition}

\begin{definition}[A discrete evaluation operator on $\Rep$]
To speak about convergence to $\dfix$ inside the time‑free space $\Rep$, we
postulate a map
\[
  \operatorname{Eval}:\Rep\to\Rep
\]
with the following properties:
\begin{enumerate}[label=(E\arabic*)]
  \item (Fixed point) $\operatorname{Eval}(\dfix)=\dfix$.
  \item (Strict decrease away from $\dfix$) For all $x\in\Rep\setminus\{\dfix\}$,
        \[
          \|\operatorname{Eval}(x)\| = \|x\|-1.
        \]
\end{enumerate}
This can be seen as a discrete analogue of a contraction step (e.g.
``delete one bit of description''), and is the only additional ingredient
needed to make the limit claim fully formal.
\end{definition}

%%%%%%%%%%%%%%%%%%%%%%%%%%%%%%%%%%%%%%%%%%%%%%%%%%%%%%%%%%%%
\section{Formal statement that a proof ``exists''}
%%%%%%%%%%%%%%%%%%%%%%%%%%%%%%%%%%%%%%%%%%%%%%%%%%%%%%%%%%%%

\begin{theorem}[Catuskoti $0$d$‑$isomorphism]\label{thm:c4-dfix}
Let $x\in\Rep$. Define the sequence in the static space
\[
  v_0:=x,\qquad v_{k+1}:=\operatorname{Eval}(v_k)\quad(k\ge 0).
\]
Then $\{\|v_k\|\}_{k\ge0}$ is a strictly decreasing sequence of non‑negative
integers until it reaches $0$, hence there exists $N\in\mathbb{N}$ such that
\[
  v_N=\dfix,\qquad\text{and therefore }\lim_{k\to\infty}v_k=\dfix.
\]
Moreover, the categorical maps $\iota:\Rep\hookrightarrow\Prf$ and
$\phi:\Prf\twoheadrightarrow\Rep$ provide the time‑axis bookkeeping:
$\iota$ adds one unit of time and $\phi$ removes it, satisfying
$\|\phi(p)\|=\|p\|-1$ for all $p\in\Prf$.
\end{theorem}
\begin{proof}
By (E2), if $v_k\neq\dfix$ then $\|v_{k+1}\|=\|v_k\|-1$; by (E1) the sequence
stays at $\dfix$ once reached.
Since $\|v_0\|\in\mathbb{N}$ and decreases by $1$ at each step until it hits
$0$, after $N=\|v_0\|$ steps we obtain $\|v_N\|=0$.
By definition of $\dfix$ as the unique element of distance $0$, we have
$v_N=\dfix$.
The statement about $\phi$ is exactly the contraction property proved from
(D1)--(D2).
\end{proof}

%%%%%%%%%%%%%%%%%%%%%%%%%%%%%%%%%%%%%%%%%%%%%%%%%%%%%%%%%%%%
\section{Summary}
%%%%%%%%%%%%%%%%%%%%%%%%%%%%%%%%%%%%%%%%%%%%%%%%%%%%%%%%%%%%

\begin{itemize}[nosep]
  \item The act of proving is captured categorically by a
        \textbf{monomorphism} $\iota$ (information‑preserving embedding) and an
        \textbf{epimorphism} $\phi$ (information‑recovering projection).
  \item Both maps are \textbf{split}, guaranteeing exact recovery of the
        original static data.
  \item The distance $\|X\|=K(X)+\ell(X)$ provides a quantitative measure of
        ``information + time''.  The evaluation map $\phi$ is a
        \textbf{contraction}, ensuring monotone decrease of this distance.
  \item Iterating $\phi\!\circ\!\iota$ yields a strictly decreasing integer
        sequence that inevitably reaches the unique $0$‑dimensional fixed
        point $\dfix$.
  \item Hence the statement “a proof exists for $(s_x,s_y)$’’ is precisely the
        existence of the morphisms $\iota$ and $\phi$ satisfying the
        properties above.
\end{itemize}

%%%%%%%%%%%%%%%%%%%%%%%%%%%%%%%%%%%%%%%%%%%%%%%%%%%%%%%%%%%%
\section{Further Remarks}
%%%%%%%%%%%%%%%%%%%%%%%%%%%%%%%%%%%%%%%%%%%%%%%%%%%%%%%%%%%%

\begin{enumerate}
  \item The categorical picture makes the previously informal phrase
        “adding a time‑axis and then evaluating’’ mathematically precise.
  \item The discrete evaluation operator $\operatorname{Eval}$ is the analogue
        of a Banach contraction step in this setting; it guarantees
        convergence to the unique $0$‑dimensional fixed point $\dfix$.
  \item One may enrich the framework by specifying concrete choices for
        $K$ (e.g. literal string length) and for the set of time tokens
        $\Tset$ (e.g. a singleton $\{*\}$), which turns the abstract
        construction into an implementable algorithm.
\end{enumerate}

\end{document}
